% The histogram itself: what the instrument produces, and how to read its shape.
% The bar heights are illustrative of the shape to expect, not measured data,
% which is why the caption says so and the axis carries no counts.
\begin{tikzpicture}[font=\sffamily\scriptsize, x=1cm, y=1cm]

% ---------------- axes
\draw[taxis] (-4.6,0) -- (6.2,0);
\draw[taxis] (-4.6,0) -- (-4.6,3.4);
\node[signame, anchor=south, rotate=90] at (-5.1,1.7) {samples per bin};
\node[anchor=north] at (0.8,-0.62) {period minus 280,000 cycles, in nanoseconds. One bin is 8 cycles, 28.6 ns};

% ---------------- the bars: one bin is 8 cycles, which is 28.6 ns
\foreach \x/\h in {%
  -1.30/0.04, -1.10/0.07, -0.92/0.13, -0.73/0.30, -0.55/0.62, -0.37/1.15,
  -0.18/1.95,  0.00/2.90,  0.18/2.15,  0.37/1.30,  0.55/0.70,  0.73/0.34,
   0.92/0.19,  1.10/0.12,  1.28/0.09,  1.47/0.07,  1.65/0.06,  1.83/0.05,
   2.02/0.05,  2.20/0.04,  2.38/0.04,  2.57/0.03,  2.75/0.03,  2.93/0.03,
   3.12/0.02,  3.30/0.02,  3.67/0.02,  4.03/0.02,  4.40/0.015, 4.77/0.015,
   5.13/0.01,  5.50/0.01}
  { \fill[kercol, draw=linecol, line width=0.3pt] (\x-0.08,0) rectangle (\x+0.08,\h); }

% ---------------- the three numbers that get reported
\draw[deadline] (0.02,0) -- (0.02,3.2);
\node[tlabel, anchor=south west] at (0.06,3.0) {median};
\draw[deadline] (3.30,0) -- (3.30,2.4);
\node[tlabel, anchor=south] at (3.30,2.4) {99.9th percentile};
\draw[deadline] (5.50,0) -- (5.50,1.5);
\node[tlabel, anchor=south] at (5.50,1.5) {maximum, the number to quote};

% ---------------- what is deliberately not reported
\node[umlnote, text width=30mm, anchor=north west] at (-4.5,3.3)
  {The mean is not reported at all. A control loop is not harmed by its average
   period and it is harmed by its worst one.};

% ---------------- the axis ticks, in nanoseconds
\foreach \x/\lab in {-1.30/-180, -0.55/-80, 0/0, 0.55/80, 1.47/210, 3.30/460, 5.50/770}
  { \draw[tick] (\x,-0.12) -- (\x,3.2); \node[tlabel, anchor=north] at (\x,-0.16) {\lab}; }

% ---------------- how to read the shape
\node[umlnote, text width=33mm, anchor=north west] at (-4.6,-1.15)
  {\textbf{One narrow peak.} The timer is doing its job and the release path is
   quiet. This is the shape to expect on the handler route.};
\node[umlnote, text width=33mm, anchor=north] at (0.9,-1.15)
  {\textbf{Two peaks.} Something else takes the processor at a regular rate.
   Look for another periodic handler, or a kernel tick that is not this
   chapter's timer.};
\node[umlnote, text width=33mm, anchor=north east] at (6.2,-1.15)
  {\textbf{A long tail to the right.} A rare handler, a critical section, or a
   cache effect. The tail is why the maximum over an hour is the number and the
   first reading is not.};

\node[chlabel, anchor=north, text width=105mm, align=center] at (0.8,-3.95)
  {Anything outside the 64 bins is counted separately rather than clamped into
   the end bins, so a rare large excursion cannot hide inside the distribution.};
\end{tikzpicture}
